\section{Source and common interfaces}
This application concerns the complete mathematical result set of
J.~D. Key and B.~G. Rodrigues, \emph{LCD codes from adjacency matrices of
graphs}, author manuscript dated 4 August 2017 (F017, seventeen pages).
The source's constructions, classifications, finite tables and recorded
questions remain distinguished. Several statements as printed need
correction; the qualifications below concern this exact edition and do
not attribute its errors to earlier papers without examining them.

Throughout, a graph is finite, simple and undirected, with at least one
vertex. Over $\mathbb F_p$, let $A$ be its adjacency matrix and set
\[
 B=J-I-A,\qquad R=I+A,\qquad \bar R=I+B=J-A,
 \quad C=\im A,\quad\bar C=\im B,\quad RC=\im R.
\]
An orthogonal complement is for the usual coordinate form. A code is
LCD if its intersection with its orthogonal complement is zero.
Following the source, $C$ is RLCD if $C=(RC)^\perp$; for a regular graph
it is CLCD if $C=\bar C^\perp$. These stronger properties should not be
confused with LCD itself.

The complete common paper supplies the orthogonal-projector theorem,
its exact constant-line defect, the scalar intersection-count criterion,
the simplex boundary argument, the ternary inclusion argument and the
transitive systematic-message certificate. The exact dependency map
binds their previously reviewed proof texts and required computations.
They are actual prerequisites, not free references to an uncharged proof
bank. This application supplies the source-specific substitutions,
complete output tables and qualifications. The common paper and all
required finite certificates must accompany it in a complete edition.
The classical spectra and some distances explicitly quoted in the
source from earlier work remain separately attributed inputs.
